\section{The Learned Encoder}\label{supp:learned}

This section gives the details of the encoder network of Section~\ref{sec:learned-encoder} and of its two evaluations in Section~\ref{sec:learned}. It describes how the network is built and trained and how its conversion was developed. It then reports how the network performs as a classifier of its own under the nested cross-validation of the main benchmark, and what it changes inside ArrowFlow.

\subsection{Network, Conversion and Training}\label{supp:learned-method}

The encoder network is the tensor layer of the ArrowFlow code (the class \texttt{TensorNet} in \texttt{arrowflow/arrowflow.py}). It has three fully connected layers of width $e$. Instance normalization and a rectified linear unit follow each of the first two, and the weights of those two layers start uniform on $[0,1)$, as the code draws them. The instance normalization has no learned scale or shift. It rescales the $e$ values of each example to mean zero and unit variance. The network's input is an example's features after the imputation and standardization of Section~\ref{sec:encoder}, without the polynomial expansion, and the code negates this input by convention. Its output is the score vector $h\in\R^e$, and the input ranking is $\pi_x=\orderdesc(h)$.

The network is trained together with one ranking layer of $C$ class filters, the output layer of Section~\ref{sec:training}. Both use the same gate rule. A misclassified example is always accepted, and a correctly classified one with probability $p_c=0.1$. The class filters learn by the output layer's rule at learning rate $\eta=0.1$. For the accepted examples of a batch of 32, the target conversion of Equation~\eqref{eq:target} gives the target scores $\tilde h$. The target is held fixed during the step. The network then takes one step of Adam \citep{kingma2015adam}, a standard method of gradient-based training, on $\tfrac12\lVert h-\tilde h\rVert^2$, averaged over those examples, at learning rate $0.01$ (Section~\ref{supp:learned-nested} also tries $0.003$). The code's step schedule multiplies this learning rate by $0.7$ after 20, 40, 50 and 100 updates. So after the first 100 updates it stays at $0.7^4\approx0.24$ of its starting value. The schedule takes these steps from the code's default of 200 updates, whatever the length of training. Training runs for 60 passes over the training rows in random batches, with no validation checkpoint. That is $60\lceil n_{\mathrm{tr}}/32\rceil$ updates for $n_{\mathrm{tr}}$ training rows, each on 32 rows drawn at random. The trained network and its layer of class filters form a classifier of their own. Section~\ref{supp:learned-nested} reads it in two ways: by its class filters, as the output layer does, or by a nearest-neighbor vote on the encoder's ranking, as ArrowFlow does. Inside ArrowFlow, the layer of class filters is set aside, and the trained network is the view's encoder.

The class filters and the encoder network draw their acceptances separately, so the two gates can disagree about a correctly classified example. Suppose the encoder network accepts it and the class filters' vote rejects it. Then the example has no motion toward its true filter, so $m_y=0$. Its target only moves it away from the nearest wrong filter, the filter of another class nearest to $\pi_x$ in footrule distance. In the reverse case the true class filter receives its vote, and the network receives no target from the example. The away term reads the class filters after this batch's reorder, and $m_y$ is computed before it. ArrowFlow's own output vote has no away term, because it only attracts.

All of this runs on the ArrowFlow code itself. A module, \nolinkurl{experiments/tensor_rank_sandbox/core_hybrid.py}, builds the code's network of one tensor layer and one ranking layer. On that network only, it replaces the method that turns the ranking layer's motions into a tensor update. The file \texttt{arrowflow.py} is unchanged. Where a single linear layer replaces the encoder network, as in Sections~\ref{supp:learned-development} and~\ref{supp:learned-nested}, its weights start as standard normal draws divided by $\sqrt d$. Its biases start as standard normal draws times $0.1$.

\subsection{How the Conversion Was Developed}\label{supp:learned-development}

The code's own conversion for such networks, the method that the module replaces, has two defects. Its motion $m$ is a coordinate's position in the input ranking minus its position in the filter, the motion of Equation~\eqref{eq:motion} with its sign reversed. It forms the target $h-\epsilon\,m\,\lvert h\rvert$ and takes an $L_1$ loss against it. But it does not detach the target from the network, so the gradient flows through the target as well. The loss is therefore $\epsilon\lvert m\rvert\lvert h\rvert$, and its gradient, $\epsilon\lvert m\rvert\operatorname{sign}(h)$, is the same whatever the signs of the motions. It only pulls the scores of displaced coordinates toward zero. Even detached, the formula would step each score by $-\operatorname{sign}(m)$. Because the code orders the scores decreasingly, that step moves a coordinate away from the position its filter asks for. Numerically, the gradient of the code's conversion does not change when every motion changes sign, and it equals $\epsilon\lvert m\rvert\operatorname{sign}(h)$ divided by the number of scores. We applied five updates of each conversion to one batch of Digits, from the same start and with the class filters held fixed. The check accepted every example and used no away term. The code's conversion moved the rankings of a linear encoder from a mean footrule distance of 1381 to their true class filters to 922, through the pull toward zero alone. The detached formula pushed them out to 1723, and the target conversion brought them to 771 (\texttt{conversion\_check.py}).

The target conversion was designed in a sandbox, a small separate program that reimplements the output layer's rule in NumPy (\nolinkurl{experiments/tensor_rank_sandbox/sandbox.py}). The sandbox used a linear encoder and chose its settings by cross-validation on the 70\% development portions of Iris, Wine, Breast cancer and Digits. It was compared with two other conversions: the sign of the motion alone, and the motion scaled by the mean gap between sorted scores. The search ran in three stages. The first tried five learning rates for each conversion, without the away term. There, each at its best learning rate, the sign of the motion reached the highest mean development accuracy, 92.4\%, against 92.2\% for the scaled motion and 91.2\% for the target conversion. Only the target conversion was carried forward, and the later stages tuned it alone. They varied the learning rate and the vocabulary size, and tried the away term, a validation checkpoint, a decaying learning rate and a frozen layer of class filters. The configuration with the highest mean development accuracy was fixed before any held-out score was read. It has a vocabulary of 64, the away term at one half and a learning rate of 0.01 that decays linearly over 150 passes. It was then tested with five seeds on the held-out 30\% of those four datasets and on 70/30 splits of eight more datasets of this paper: Vehicle, Segment, balance-scale, mfeat-zernike, ionosphere, diabetes, vertebra-column and steel-plates-fault. The trained linear encoder was more accurate than the frozen one on all 12 datasets. The same signal with its coordinates shuffled was less accurate than the frozen encoder on 11 of the 12.

The target conversion was then ported into the ArrowFlow code. The port kept the vocabulary of 64, the away term at one half and the learning rate of 0.01. It trained with Adam and the code's step schedule for 60 passes, in place of plain gradient steps with a linear decay over 150 passes. Run this way on the same twelve splits, the code's own encoder network reached a mean accuracy of 86.5\% with the target conversion. It reached 70.5\% when left at its random start, 73.4\% with the code's own conversion, 87.2\% with that conversion detached and its sign reversed, and 55.7\% with the shuffled signal. The target conversion, fixed in the sandbox, was kept. It trails the repaired code conversion by 0.7 points with the code's encoder network, and it leads by 0.7 points with a linear encoder, 84.2\% against 83.5\%. The MLP comparison model of this paper, tuned in the same way, reached 90.3\%. These single splits cover twelve of the seventeen datasets of this paper. Their test sets differ from the test sets of the main benchmark. Every choice made on them came before the runs of Sections~\ref{supp:learned-nested} and~\ref{supp:learned-swap}: the conversion with its settings, the code's encoder network and the grids of those sections. Those runs kept the settings of the port, with two exceptions. The nested study also tried a vocabulary of 128 and a learning rate of 0.003. The encoder swap took its vocabulary from ArrowFlow's configuration. Apart from the comparisons of this subsection, no result in this paper uses the code's own conversion.

\subsection{The Classifier under Nested Cross-Validation}\label{supp:learned-nested}

This study asks whether the converted motions train the encoder network at all. It uses the data and the 15 test sets of the main benchmark. Like every model of Table~\ref{tab:main}, each of its models is tuned on the inner folds and refitted with the three fitting seeds. Its design was fixed in the header of \nolinkurl{experiments/tensor_rank_sandbox/nested.py} and committed before the run (Table~\ref{tab:s-protocols}). Three versions of the classifier of Section~\ref{supp:learned-method} are compared. The first is the trained encoder network, with vocabulary size $e\in\{64,128\}$ and learning rate $0.003$ or $0.01$ chosen on the inner folds. The second is the same network never trained, with $e$ chosen likewise, and the third a trained single linear layer in place of the network, on the grid of the first. The inner folds choose by the accuracy of the class filters' readout at the first fitting seed. Unlike the stochastic models of Table~\ref{tab:main}, this study does not rescore its three best candidates with all three seeds. Each fit is also read by a nearest-neighbor vote on the encoder's ranking, with its neighbor count and weighting chosen on the training rows, as in ArrowFlow. Three contrasts were fixed before the run, each paired over the 15 test sets with the corrected resampled $t$ and a Holm adjustment across the seventeen datasets. Two read the classifier by its class filters: the trained minus the untrained encoder, and the trained encoder minus the MLP. The third is the trained encoder read by the nearest-neighbor vote minus ArrowFlow.

% Rendered by render_tables.py from run files; do not edit by hand.
\begin{table}[!htbp]
\caption{\textbf{The learned encoder under nested cross-validation: error.} Mean error (\%) over the 15 test sets of the main benchmark. The first four columns are the classifier of Section~\ref{supp:learned-method}, an encoder network followed by one ranking layer of class filters. It is read by the class filters or by a nearest-neighbor vote on the encoder's own ranking, and its encoder is the trained network, the same network never trained, or a trained single linear layer. Each is tuned on the inner folds at the first fitting seed and refitted with the three fitting seeds. Unlike the stochastic models of the main benchmark, it does not rescore its three best candidates with all three seeds. The untrained encoder keeps its initial weights, uniform on $[0,1)$ in its first two layers. ArrowFlow and the MLP are the registered results on the same test sets. HCV: hepatitis C virus.}
\label{tab:s-learned-nested}
\centering
\scriptsize
\setlength{\tabcolsep}{3pt}
\resizebox{\ifdim\width>\linewidth\linewidth\else\width\fi}{!}{%
\begin{tabular}{lrrrrrr}
\toprule
Dataset & \begin{tabular}[b]{@{}r@{}}Trained encoder,\\class filters\end{tabular} & \begin{tabular}[b]{@{}r@{}}Untrained encoder,\\class filters\end{tabular} & \begin{tabular}[b]{@{}r@{}}Trained linear enc.,\\class filters\end{tabular} & \begin{tabular}[b]{@{}r@{}}Trained encoder,\\kNN readout\end{tabular} & ArrowFlow & MLP \\
\midrule
\multicolumn{7}{l}{\emph{The seven benchmark datasets}} \\
Iris & 5.0 & 7.0 & 11.1 & 4.6 & 5.0 & 4.1 \\
Wine & 1.9 & 5.2 & 3.3 & 1.9 & 3.2 & 1.9 \\
Breast cancer & 2.8 & 5.3 & 3.2 & 2.8 & 2.7 & 2.9 \\
Wine quality & 43.4 & 45.8 & 43.6 & 30.1 & 27.7 & 32.3 \\
Vehicle & 25.0 & 48.7 & 24.0 & 24.4 & 19.9 & 17.2 \\
Segment & 6.3 & 17.9 & 9.6 & 3.4 & 2.6 & 2.8 \\
Digits & 3.0 & 20.9 & 6.5 & 2.4 & 2.8 & 2.3 \\
\midrule
\multicolumn{7}{l}{\emph{The ten further datasets}} \\
balance-scale & 3.7 & 16.3 & 10.3 & 2.5 & 6.0 & 2.2 \\
mfeat-zernike & 17.0 & 34.3 & 20.4 & 17.2 & 19.3 & 18.9 \\
ionosphere & 8.0 & 18.0 & 13.2 & 7.9 & 11.2 & 8.0 \\
vertebra-column & 18.6 & 21.8 & 17.0 & 17.8 & 16.9 & 15.4 \\
diabetes & 29.3 & 30.5 & 26.4 & 26.3 & 24.0 & 23.5 \\
banknote-authentication & 0.0 & 4.0 & 3.0 & 0.0 & 0.2 & 0.0 \\
qsar-biodeg & 13.7 & 21.9 & 16.4 & 13.4 & 12.6 & 12.0 \\
steel-plates-fault & 34.7 & 44.1 & 38.2 & 29.4 & 24.8 & 24.2 \\
climate-model-simulation-crashes & 5.9 & 17.7 & 6.7 & 5.8 & 8.3 & 5.6 \\
HCV & 74.4 & 74.9 & 75.3 & 74.9 & 75.6 & 75.0 \\
\midrule
\multicolumn{7}{l}{\emph{All datasets}} \\
Mean over the seventeen & 17.2 & 25.5 & 19.3 & 15.6 & 15.5 & 14.6 \\
\bottomrule
\end{tabular}
}
\end{table}

% Rendered by render_tables.py from run files; do not edit by hand.
\begin{table}[!htbp]
\caption{\textbf{The learned encoder under nested cross-validation: the three contrasts fixed before the run.} Accuracy differences, in percentage points, with the 95\% corrected resampled-$t$ interval (test-to-training ratio 0.25, 14 degrees of freedom); positive favors the first-named model, and each contrast has its own Holm adjustment across the seventeen datasets. The first two contrasts read the classifier of Section~\ref{supp:learned-method} by its class filters; the third reads its trained encoder's ranking with a nearest-neighbor vote, the readout ArrowFlow uses. The trained encoder is more accurate than the untrained one on all 17 datasets, significantly on 11. Against the MLP it is ahead on 3, significantly on none, and significantly behind on four; with the nearest-neighbor readout it is ahead of ArrowFlow on 9, significantly on one, and significantly behind on two. HCV: hepatitis C virus.}
\label{tab:s-learned-nested-contrasts}
\centering
\scriptsize
\setlength{\tabcolsep}{3pt}
\resizebox{\ifdim\width>\linewidth\linewidth\else\width\fi}{!}{%
\begin{tabular}{lrrrrrr}
\toprule
Dataset & \begin{tabular}[b]{@{}r@{}}Trained $-$ untrained\\encoder (pp)\end{tabular} & Holm $p$ & \begin{tabular}[b]{@{}r@{}}Trained encoder\\$-$ MLP (pp)\end{tabular} & Holm $p$ & \begin{tabular}[b]{@{}r@{}}kNN readout $-$\\ArrowFlow (pp)\end{tabular} & Holm $p$ \\
\midrule
\multicolumn{7}{l}{\emph{The seven benchmark datasets}} \\
Iris & +2.07 [$-$1.75, +5.90] & 1.000 & $-$0.81 [$-$3.11, +1.48] & 1.000 & +0.44 [$-$1.75, +2.64] & 1.000 \\
Wine & +3.31 [+0.95, +5.66] & 0.056 & $-$0.07 [$-$1.38, +1.25] & 1.000 & +1.31 [$-$1.06, +3.67] & 1.000 \\
Breast cancer & +2.48 [+0.96, +4.00] & 0.025 & +0.06 [$-$0.89, +1.01] & 1.000 & $-$0.12 [$-$1.22, +0.99] & 1.000 \\
Wine quality & +2.43 [$-$2.13, +6.99] & 1.000 & $-$11.05 [$-$14.92, $-$7.17] & $<$0.001 & $-$2.42 [$-$4.94, +0.09] & 0.630 \\
Vehicle & +23.72 [+21.27, +26.17] & $<$0.001 & $-$7.79 [$-$10.97, $-$4.61] & 0.002 & $-$4.48 [$-$7.12, $-$1.84] & 0.040 \\
Segment & +11.65 [+10.03, +13.28] & $<$0.001 & $-$3.49 [$-$4.26, $-$2.71] & $<$0.001 & $-$0.83 [$-$1.54, $-$0.11] & 0.316 \\
Digits & +17.89 [+16.57, +19.22] & $<$0.001 & $-$0.71 [$-$1.55, +0.14] & 1.000 & +0.35 [$-$0.27, +0.97] & 1.000 \\
\midrule
\multicolumn{7}{l}{\emph{The ten further datasets}} \\
balance-scale & +12.59 [+10.39, +14.78] & $<$0.001 & $-$1.55 [$-$3.61, +0.52] & 1.000 & +3.52 [+2.26, +4.78] & $<$0.001 \\
mfeat-zernike & +17.27 [+14.83, +19.72] & $<$0.001 & +1.84 [$-$0.35, +4.04] & 1.000 & +2.12 [+0.45, +3.79] & 0.213 \\
ionosphere & +10.00 [+6.81, +13.20] & $<$0.001 & $-$0.00 [$-$2.41, +2.40] & 1.000 & +3.29 [$-$0.52, +7.10] & 0.853 \\
vertebra-column & +3.23 [+0.27, +6.18] & 0.172 & $-$3.12 [$-$7.03, +0.79] & 1.000 & $-$0.90 [$-$4.46, +2.67] & 1.000 \\
diabetes & +1.17 [$-$4.51, +6.85] & 1.000 & $-$5.79 [$-$9.71, $-$1.86] & 0.090 & $-$2.29 [$-$5.29, +0.71] & 1.000 \\
banknote-authentication & +3.93 [+3.14, +4.71] & $<$0.001 & $-$0.04 [$-$0.23, +0.15] & 1.000 & +0.15 [$-$0.23, +0.52] & 1.000 \\
qsar-biodeg & +8.12 [+6.11, +10.13] & $<$0.001 & $-$1.70 [$-$3.73, +0.34] & 1.000 & $-$0.80 [$-$2.64, +1.03] & 1.000 \\
steel-plates-fault & +9.38 [+7.02, +11.74] & $<$0.001 & $-$10.49 [$-$12.96, $-$8.02] & $<$0.001 & $-$4.58 [$-$6.02, $-$3.14] & $<$0.001 \\
climate-model-simulation-crashes & +11.77 [+8.24, +15.30] & $<$0.001 & $-$0.29 [$-$1.57, +0.99] & 1.000 & +2.55 [+0.71, +4.39] & 0.140 \\
HCV & +0.48 [$-$1.99, +2.95] & 1.000 & +0.57 [$-$1.71, +2.85] & 1.000 & +0.75 [$-$1.99, +3.48] & 1.000 \\
\bottomrule
\end{tabular}
}
\end{table}


The trained encoder is more accurate than the untrained one on all 17 datasets, significantly on 11 (Table~\ref{tab:s-learned-nested-contrasts}), and its mean error over the datasets is 17.2\%, against 25.5\% untrained. Read by its class filters, the classifier trails the MLP. It is ahead on 3 datasets, significantly on none, and significantly behind on 4, Wine quality, Vehicle, Segment and steel-plates-fault. Read by a nearest-neighbor vote on the trained encoder's ranking, from a single view with no hidden ranking layer, its mean error of 15.6\% is close to ArrowFlow's 15.5\%. It is ahead of ArrowFlow on 9 datasets, significantly on balance-scale, and significantly behind on Vehicle and steel-plates-fault. The trained single linear layer reaches a mean error of 19.3\%, between the trained and the untrained network.

\subsection{The Encoder Swap}\label{supp:learned-swap}

This study asks whether ArrowFlow gains from the trained encoder, with nothing else retuned. For every outer split, ArrowFlow keeps the configuration that its inner-fold tuning chose on that split's training partition. Its encoded vocabulary size, degree and augmentation switch follow from that choice and that partition, as in the component ablation. Only each view's fixed encoder is replaced, by an encoder network. The encoder network has the settings of Section~\ref{supp:learned-method}, with $e$ equal to the encoded vocabulary size of the selected configuration. Its width therefore ranges from 16 to 128 over the splits, while the nested study tried only 64 and 128. It is trained on the view's training rows with the view's seed. It reads the features without the polynomial expansion, so the degree of the configuration does not apply. Without projection strategies, no view is target-aware or calibrated, and the views differ only by their seeds. The untrained arm keeps the same network at its random start, and the fixed arm is ArrowFlow's registered predictions.

The design was fixed in the header of \nolinkurl{experiments/tensor_rank_sandbox/encoder_swap.py} and committed before the run (Table~\ref{tab:s-protocols}). This was after the outer results of Section~\ref{supp:learned-nested} were known, including the significant gain of the trained encoder's nearest-neighbor readout over ArrowFlow on balance-scale. The design includes the rule that reads the contrast. The trained encoder improves ArrowFlow if the mean accuracy difference over the seventeen datasets is positive and more datasets are significantly higher than significantly lower. In the reverse case the rule reads that the trained encoder worsens ArrowFlow, and otherwise that no difference is shown. The design also declared four secondary contrasts, all descriptive, each with its own Holm adjustment across the seventeen datasets. Three compare the arms with each other and with the MLP (Table~\ref{tab:s-learned-swap-controls}). The fourth sets ArrowFlow with the trained encoder against the trained encoder of Section~\ref{supp:learned-nested} read alone by its nearest-neighbor vote.

Before the run, the module refitted the fixed arm on 4 datasets at all three fitting seeds and reproduced the stored predictions exactly. None of the 12 fits differed on any test row. It used one test set of each of Iris, Segment, balance-scale and mfeat-zernike. During the run, 5 jobs were computed twice, once by each of two worker pools, and their records were identical.

% Rendered by render_tables.py from run files; do not edit by hand.
\begin{table}[!htbp]
\caption{\textbf{The encoder swap: ArrowFlow with the trained encoder.} Mean error (\%) over the 15 test sets, and the contrast fixed before the run, in percentage points: ArrowFlow with the trained encoder minus ArrowFlow with its fixed encoder. The contrast has the 95\% corrected resampled-$t$ interval (test-to-training ratio 0.25, 14 degrees of freedom) and a Holm adjustment across the seventeen datasets. For every outer split, ArrowFlow keeps the configuration that its inner-fold tuning chose on that split's training partition, and nothing is retuned. Each view's encoder is the fixed encoder of Section~\ref{sec:encoder} (the registered results), the trained encoder network, or the same network never trained. An encoder network reads the features without the polynomial expansion and uses no projection strategy, so its views differ only by their seeds. The MLP is the registered comparison model. With the trained encoder, ArrowFlow has the higher mean accuracy on 11 of the 17 datasets, significantly on one and significantly lower on none. HCV: hepatitis C virus.}
\label{tab:s-learned-swap}
\centering
\scriptsize
\setlength{\tabcolsep}{3pt}
\resizebox{\ifdim\width>\linewidth\linewidth\else\width\fi}{!}{%
\begin{tabular}{lrrrrrrr}
\toprule
Dataset & \begin{tabular}[b]{@{}r@{}}Fixed encoder\\(ArrowFlow)\end{tabular} & \begin{tabular}[b]{@{}r@{}}Trained\\encoder\end{tabular} & \begin{tabular}[b]{@{}r@{}}Untrained\\encoder\end{tabular} & MLP & \begin{tabular}[b]{@{}r@{}}Trained $-$ fixed\\Diff. (pp) [interval]\end{tabular} & \begin{tabular}[b]{@{}r@{}}Trained $-$ fixed\\$p$\end{tabular} & \begin{tabular}[b]{@{}r@{}}Trained $-$ fixed\\Holm $p$\end{tabular} \\
\midrule
\multicolumn{8}{l}{\emph{The seven benchmark datasets}} \\
Iris & 5.0 & 4.1 & 6.5 & 4.1 & +0.89 [$-$1.66, +3.43] & 0.466 & 1.000 \\
Wine & 3.2 & 1.9 & 4.0 & 1.9 & +1.31 [$-$1.31, +3.92] & 0.302 & 1.000 \\
Breast cancer & 2.7 & 2.7 & 4.4 & 2.9 & $-$0.06 [$-$1.28, +1.17] & 0.920 & 1.000 \\
Wine quality & 27.7 & 27.7 & 27.9 & 32.3 & $-$0.01 [$-$1.97, +1.94] & 0.987 & 1.000 \\
Vehicle & 19.9 & 23.6 & 30.0 & 17.2 & $-$3.71 [$-$7.07, $-$0.36] & 0.032 & 0.452 \\
Segment & 2.6 & 3.2 & 3.4 & 2.8 & $-$0.58 [$-$1.17, +0.01] & 0.055 & 0.711 \\
Digits & 2.8 & 2.3 & 4.5 & 2.3 & +0.47 [$-$0.18, +1.11] & 0.141 & 1.000 \\
\midrule
\multicolumn{8}{l}{\emph{The ten further datasets}} \\
balance-scale & 6.0 & 1.6 & 10.6 & 2.2 & +4.46 [+3.09, +5.84] & $<$0.001 & $<$0.001 \\
mfeat-zernike & 19.3 & 17.5 & 21.7 & 18.9 & +1.75 [+0.61, +2.89] & 0.005 & 0.087 \\
ionosphere & 11.2 & 7.2 & 13.5 & 8.0 & +3.99 [$-$0.20, +8.18] & 0.060 & 0.722 \\
vertebra-column & 16.9 & 16.2 & 21.8 & 15.4 & +0.72 [$-$1.67, +3.10] & 0.530 & 1.000 \\
diabetes & 24.0 & 25.7 & 26.7 & 23.5 & $-$1.69 [$-$4.27, +0.88] & 0.180 & 1.000 \\
banknote-authentication & 0.2 & 0.0 & 0.2 & 0.0 & +0.19 [$-$0.37, +0.74] & 0.484 & 1.000 \\
qsar-biodeg & 12.6 & 12.2 & 13.0 & 12.0 & +0.33 [$-$1.68, +2.33] & 0.732 & 1.000 \\
steel-plates-fault & 24.8 & 26.5 & 25.2 & 24.2 & $-$1.65 [$-$3.52, +0.23] & 0.080 & 0.883 \\
climate-model-simulation-crashes & 8.3 & 5.5 & 8.5 & 5.6 & +2.86 [+0.76, +4.96] & 0.011 & 0.165 \\
HCV & 75.6 & 75.4 & 74.7 & 75.0 & +0.20 [$-$2.63, +3.03] & 0.881 & 1.000 \\
\midrule
\multicolumn{8}{l}{\emph{All datasets}} \\
Mean over the seventeen & 15.5 & 14.9 & 17.5 & 14.6 & +0.56 & -- & -- \\
\bottomrule
\end{tabular}
}
\end{table}

% Rendered by render_tables.py from run files; do not edit by hand.
\begin{table}[!htbp]
\caption{\textbf{The encoder swap: controls.} Accuracy differences between the arms of Table~\ref{tab:s-learned-swap}, in percentage points, with the 95\% corrected resampled-$t$ interval (test-to-training ratio 0.25, 14 degrees of freedom); positive favors the first-named arm. The design declared these contrasts secondary and descriptive, and only the contrast of Table~\ref{tab:s-learned-swap} primary. Each has its own Holm adjustment across the seventeen datasets, and a Holm $p$ below 0.05 counts as significant within it. The trained encoder beats the untrained one on 15 of the 17 datasets, significantly on four. The untrained encoder is behind the fixed one on 16, significantly on three. HCV: hepatitis C virus.}
\label{tab:s-learned-swap-controls}
\centering
\scriptsize
\setlength{\tabcolsep}{3pt}
\resizebox{\ifdim\width>\linewidth\linewidth\else\width\fi}{!}{%
\begin{tabular}{lrrrrrr}
\toprule
Dataset & \begin{tabular}[b]{@{}r@{}}Trained $-$ untrained\\(pp) [interval]\end{tabular} & Holm $p$ & \begin{tabular}[b]{@{}r@{}}Untrained $-$ fixed\\(pp) [interval]\end{tabular} & Holm $p$ & \begin{tabular}[b]{@{}r@{}}Trained $-$ MLP\\(pp) [interval]\end{tabular} & Holm $p$ \\
\midrule
\multicolumn{7}{l}{\emph{The seven benchmark datasets}} \\
Iris & +2.37 [$-$1.55, +6.29] & 1.000 & $-$1.48 [$-$4.63, +1.67] & 1.000 & +0.00 [$-$1.60, +1.60] & 1.000 \\
Wine & +2.11 [$-$1.01, +5.24] & 1.000 & $-$0.81 [$-$3.70, +2.09] & 1.000 & $-$0.00 [$-$1.54, +1.53] & 1.000 \\
Breast cancer & +1.66 [+0.12, +3.20] & 0.368 & $-$1.72 [$-$3.07, $-$0.37] & 0.198 & +0.12 [$-$0.48, +0.71] & 1.000 \\
Wine quality & +0.24 [$-$1.42, +1.90] & 1.000 & $-$0.26 [$-$2.25, +1.73] & 1.000 & +4.64 [+2.56, +6.72] & 0.005 \\
Vehicle & +6.41 [+3.37, +9.45] & 0.007 & $-$10.13 [$-$13.75, $-$6.50] & $<$0.001 & $-$6.42 [$-$9.62, $-$3.22] & 0.012 \\
Segment & +0.24 [$-$0.51, +0.99] & 1.000 & $-$0.82 [$-$1.35, $-$0.28] & 0.075 & $-$0.39 [$-$1.25, +0.47] & 1.000 \\
Digits & +2.20 [+0.71, +3.69] & 0.083 & $-$1.73 [$-$3.04, $-$0.43] & 0.169 & $-$0.04 [$-$0.79, +0.70] & 1.000 \\
\midrule
\multicolumn{7}{l}{\emph{The ten further datasets}} \\
balance-scale & +9.03 [+7.34, +10.72] & $<$0.001 & $-$4.57 [$-$6.14, $-$3.00] & $<$0.001 & +0.59 [$-$0.80, +1.97] & 1.000 \\
mfeat-zernike & +4.15 [+3.06, +5.24] & $<$0.001 & $-$2.40 [$-$3.66, $-$1.14] & 0.017 & +1.34 [$-$0.47, +3.15] & 1.000 \\
ionosphere & +6.24 [+3.61, +8.86] & 0.002 & $-$2.25 [$-$4.92, +0.42] & 0.927 & +0.79 [$-$1.08, +2.67] & 1.000 \\
vertebra-column & +5.63 [+1.37, +9.88] & 0.145 & $-$4.91 [$-$9.14, $-$0.68] & 0.287 & $-$0.72 [$-$3.99, +2.55] & 1.000 \\
diabetes & +0.97 [$-$0.82, +2.76] & 1.000 & $-$2.66 [$-$6.22, +0.89] & 1.000 & $-$2.21 [$-$4.19, $-$0.24] & 0.430 \\
banknote-authentication & +0.22 [$-$0.17, +0.60] & 1.000 & $-$0.03 [$-$0.31, +0.25] & 1.000 & +0.00 [+0.00, +0.00] & 1.000 \\
qsar-biodeg & +0.76 [$-$1.34, +2.86] & 1.000 & $-$0.43 [$-$1.69, +0.82] & 1.000 & $-$0.19 [$-$2.10, +1.72] & 1.000 \\
steel-plates-fault & $-$1.27 [$-$2.87, +0.34] & 1.000 & $-$0.38 [$-$2.19, +1.42] & 1.000 & $-$2.30 [$-$3.64, $-$0.97] & 0.036 \\
climate-model-simulation-crashes & +3.05 [+1.01, +5.08] & 0.083 & $-$0.19 [$-$0.60, +0.23] & 1.000 & +0.14 [$-$0.93, +1.22] & 1.000 \\
HCV & $-$0.66 [$-$3.48, +2.16] & 1.000 & +0.86 [$-$1.27, +2.99] & 1.000 & $-$0.43 [$-$3.44, +2.57] & 1.000 \\
\bottomrule
\end{tabular}
}
\end{table}


With the trained encoder, ArrowFlow has the higher mean accuracy on 11 of the 17 datasets, significantly on one and significantly lower on none. Its mean error over the datasets falls from 15.5\% to 14.9\%, against 17.5\% with the untrained network and 14.6\% for the MLP. The rule fixed before the run therefore reads that the trained encoder improves ArrowFlow. The one significant gain is on balance-scale, one of the twelve datasets on which the conversion was developed (Section~\ref{supp:learned-development}). Across the datasets the gain is not established. A post hoc Wilcoxon signed-rank test over the 17 mean differences, one-sided as in Table~\ref{tab:s-signed-rank}, gives $p=0.095$.

The trained encoder beats the untrained one on 15 of the 17 datasets, significantly on 4. The untrained encoder is behind the fixed encoder on 16, significantly on 3 (Table~\ref{tab:s-learned-swap-controls}). In the fourth secondary contrast, ArrowFlow with the trained encoder has the higher mean accuracy on 15 of the 17 datasets, by 0.67 points on average, and no difference is significant. Against the best tuned classical model on each dataset, ArrowFlow with the trained encoder is behind on 15 datasets and equal on banknote-authentication, where both are perfect. It is ahead on balance-scale, where its error of 1.6\% is the lowest of all the tuned models. This is descriptive, and its lead over the MLP there is not significant.
